% C055 — 原有结论保留，并把适用条件及切线推广写完整。
\begin{statementbox}
\textbf{\textcolor{CStatement}{条件}}\par\smallskip
设抛物线
\[
 C:\ y^2=2px\quad(p>0),\qquad F\!\left(\frac p2,0\right).
\]
一条\textbf{经过焦点 $F$ 的直线}与抛物线交于\textbf{两个不同的点}
$A(x_1,y_1)$、$B(x_2,y_2)$，则 $AB$ 称为焦点弦。

\medskip
\textbf{\textcolor{CStatement}{两条必记定值}}\par\smallskip
\[
 \boxed{\ x_1x_2=\frac{p^2}{4}\ },\qquad
 \boxed{\ y_1y_2=-p^2\ }.
\]
\begin{itemize}[leftmargin=*,itemsep=0.25em]
 \item 横坐标积：\textbf{正数}，等于焦点横坐标的平方。
 \item 纵坐标积：\textbf{负数}，两端点必分布在 $x$ 轴的上下两侧。
 \item 无论弦如何倾斜，这两个积都\textbf{不变}。
\end{itemize}

\medskip
\textbf{\textcolor{CStatement}{推广：焦点弦两端点切线互相垂直}}\par\smallskip
分别在 $A,B$ 处作抛物线的切线，切线斜率为
\[
 k_A=\frac{p}{y_1},\qquad k_B=\frac{p}{y_2},\qquad
 \boxed{k_Ak_B=-1}.
\]
两条切线交于 $T$，其交点\textbf{必在准线} $x=-\frac p2$ 上：
\[
 \boxed{\ T\!\left(-\frac p2,\frac{y_1+y_2}{2}\right)\ }.
\]

\medskip
\textbf{\textcolor{CStatement}{范围与变形}}\par\smallskip
\begin{itemize}[leftmargin=*,itemsep=0.2em]
 \item 直线 $y=0$ 虽过焦点，却只与抛物线交于顶点 $O$，\textbf{不构成两点焦点弦}；其余过焦点的直线均有两个不同交点。
 \item 对左开口 $y^2=-2px$（$p>0$），两个坐标积仍依次为 $\frac{p^2}{4}$、$-p^2$。
 \item 对上下开口 $x^2=\pm2py$（$p>0$），\textbf{交换 $x,y$}：$x_1x_2=-p^2$、$y_1y_2=\frac{p^2}{4}$。
\end{itemize}
\end{statementbox}
